\documentclass[11pt,a4paper]{article}
% =====================================================================
%  Shared preamble for the Part V lecture notes (lectures 19-22).
%
%  These four lectures sit outside Geron, so the notes ARE the primary
%  source and are examined on the same terms as the textbook parts.
%  Everything here is chosen to make that readable on paper: one column,
%  generous measure, and the same six-colour palette as the slides so a
%  student moving between the two is not re-learning what red means.
%
%  Build with pdflatex (twice, for the toc and the refs):
%      cd notes && latexmk -pdf lecture-19.tex
%  or  make -C notes
% =====================================================================
\usepackage[T1]{fontenc}
\usepackage[utf8]{inputenc}
\usepackage{newtxtext,newtxmath}      % Times-like: prints well, reads well
\usepackage[scaled=0.86]{inconsolata} % code, at a size that sits beside Times
\usepackage{microtype}
% newtx defines \Bbbk, \openbox, \proof and \endproof; amssymb and amsthm
% then redefine them and abort. Freeing the four names before those two
% packages load is the standard fix, and it costs nothing here: we use
% amsthm only for \newtheorem.
\let\Bbbk\relax
\usepackage{amsmath,amssymb}
\let\openbox\relax
\let\proof\relax
\let\endproof\relax
\usepackage{amsthm}
\usepackage{booktabs}
\usepackage{enumitem}
\usepackage{xcolor}
\usepackage{tcolorbox}
\usepackage{titlesec}
\usepackage{graphicx}
\usepackage[hidelinks]{hyperref}
\tcbuselibrary{skins,breakable}

% ---- the course palette, from assets/css/custom.css -------------------
\definecolor{aimlprimary}{HTML}{0B3D62}   % deep blue  - structure
\definecolor{aimlaccent} {HTML}{C0392B}   % brick red  - failures
\definecolor{aimlsuccess}{HTML}{14663A}   % green      - repairs
\definecolor{aimlmath}   {HTML}{6C3483}   % purple     - the derivation
\definecolor{aimlmuted}  {HTML}{4B5563}
\definecolor{aimlrule}   {HTML}{B0BCC7}
\definecolor{aimlsoft}   {HTML}{F4F7F9}

\usepackage[a4paper,margin=32mm,top=30mm,bottom=30mm]{geometry}
\setlength{\parskip}{0.45\baselineskip}
\setlength{\parindent}{0pt}
\linespread{1.06}

% ---- headings --------------------------------------------------------
\titleformat{\section}{\Large\bfseries\color{aimlprimary}}{\thesection}{0.7em}{}
\titleformat{\subsection}{\large\bfseries\color{aimlprimary}}{\thesubsection}{0.7em}{}
\titleformat{\subsubsection}{\bfseries\color{aimlmuted}}{}{0em}{}
\titlespacing*{\section}{0pt}{1.6\baselineskip}{0.5\baselineskip}
\titlespacing*{\subsection}{0pt}{1.2\baselineskip}{0.35\baselineskip}

% ---- boxes -----------------------------------------------------------
% One box per role, and the role is the colour. A student who has seen the
% slides already knows what each one means before reading a word of it.
\newtcolorbox{keybox}[1][]{
  breakable, enhanced, colback=aimlsoft, colframe=aimlprimary,
  boxrule=0.4pt, left=8pt, right=8pt, top=6pt, bottom=6pt,
  fonttitle=\bfseries\small, coltitle=white, colbacktitle=aimlprimary,
  attach boxed title to top left={yshift=-2mm, xshift=4mm},
  boxed title style={boxrule=0pt, arc=1pt}, #1}

\newtcolorbox{failbox}[1][]{
  breakable, enhanced, colback=white, colframe=aimlaccent,
  boxrule=0.9pt, left=8pt, right=8pt, top=6pt, bottom=6pt,
  fonttitle=\bfseries\small, coltitle=white, colbacktitle=aimlaccent,
  attach boxed title to top left={yshift=-2mm, xshift=4mm},
  boxed title style={boxrule=0pt, arc=1pt}, #1}

\newtcolorbox{derivbox}[1][]{
  breakable, enhanced, colback=white, colframe=aimlmath,
  boxrule=0.6pt, left=8pt, right=8pt, top=6pt, bottom=6pt,
  fonttitle=\bfseries\small, coltitle=white, colbacktitle=aimlmath,
  attach boxed title to top left={yshift=-2mm, xshift=4mm},
  boxed title style={boxrule=0pt, arc=1pt}, #1}

% An exercise. Deliberately the same shape as the notebook's `try` field:
% one change, and what should happen to the output.
\newtcolorbox{trybox}[1][]{
  breakable, enhanced, colback=white, colframe=aimlsuccess,
  boxrule=0.5pt, left=8pt, right=8pt, top=6pt, bottom=6pt,
  fonttitle=\bfseries\small, coltitle=white, colbacktitle=aimlsuccess,
  attach boxed title to top left={yshift=-2mm, xshift=4mm},
  boxed title style={boxrule=0pt, arc=1pt}, #1}

% ---- small semantics -------------------------------------------------
\newcommand{\scope}[1]{\textcolor{aimlmuted}{\small #1}}
\newcommand{\term}[1]{\textbf{#1}}
\newcommand{\code}[1]{\texttt{\small #1}}
\newcommand{\lecture}[1]{Lecture~#1}

\newtheorem{definition}{Definition}[section]
\newtheorem{proposition}[definition]{Proposition}

% ---- title block -----------------------------------------------------
% \notestitle{19}{Information retrieval: the lexical foundation}{SciFact (BEIR)}
\newcommand{\notestitle}[3]{%
  \thispagestyle{empty}
  \begin{flushleft}
    {\color{aimlmuted}\small\sffamily
     APPLICATIONS OF MACHINE LEARNING \quad$\cdot$\quad
     BSc Mathematics of Artificial Intelligence}\\[2mm]
    {\color{aimlrule}\rule{\textwidth}{0.8pt}}\\[4mm]
    {\color{aimlmuted}\large Lecture #1 \quad$\cdot$\quad Part V \quad$\cdot$\quad
     Extended notes}\\[3mm]
    {\Huge\bfseries\color{aimlprimary}\baselineskip=1.15\baselineskip #2\par}\vspace{4mm}
    {\color{aimlmuted}\large Dataset: #3}\\[3mm]
    {\color{aimlrule}\rule{\textwidth}{0.8pt}}
  \end{flushleft}
  \vspace{2mm}
  \begin{keybox}[title={Why these notes exist}]
    Lectures 19--22 sit outside G\'eron, the course's primary text. For those
    four lectures \emph{these notes are the primary source}, and they are
    examined on exactly the same terms as the textbook parts. Everything in
    the slide deck for this lecture is here, with the arguments written out at
    length rather than compressed onto a slide; the numbers are the ones the
    accompanying notebook prints, and every one of them is a measurement on
    the corpus named above rather than a figure quoted from a paper.
  \end{keybox}
  \vspace{1mm}
  {\small\tableofcontents}
  \vspace{2mm}
  \noindent{\color{aimlrule}\rule{\textwidth}{0.4pt}}
  \vspace{1mm}
}


\begin{document}
\notestitlefor{15}{IV}{Time series}{Chicago transit ridership}{Chapter 13}

\section{A new kind of data}

Everything so far assumed the rows were \term{exchangeable}: a random split was
valid because any row was as good as any other for training.

Not here. In a time series the rows are ordered and neighbouring rows are
alike --- one day's ridership tells you most of the next day's. Split such data
at random and a point's own near-future lands in the training set, so the score
you report is not the score you would get in deployment. It is optimistic by a
margin we will measure.

\subsection{The brief}

Forecast the number of people who will board rail and bus tomorrow, from twenty
years of daily boarding totals. The decisions it changes --- how many train
sets leave the depot, how many drivers are rostered, whether a spare crew is
held --- are made the evening before, which fixes the horizon at one day and
fixes when the forecast must exist: \emph{before the day it describes}.

\begin{keybox}[title={The two errors are not equal, and we choose a symmetric metric anyway}]
Forecast too high and empty vehicles run: fuel, crew hours, a line in the
budget. Measurable and survivable. Forecast too low and passengers are left on
platforms at peak: the cost is political and appears in no table you will be
shown.

We say that out loud now rather than discovering it later.
\end{keybox}

The label is free --- nobody annotates anything, because the next day's value is
the label for today's window. But the rows are not exchangeable: row 400 is not
a fresh draw from the same distribution as row 10, it is the same city eighteen
months later, and rows 400 and 401 are almost the same measurement made twice.

\subsection{Four lines of tidying, one of which matters}

The file ships with 7{,}701 rows and 7{,}639 distinct ones: sixty-two days
appear twice, identically, and nothing warns you. Reviewer question 4 from
Lecture 1 --- what was dropped? --- and here the answer is 62, which you only
know because you counted.

One categorical column, \code{day\_type}: 5{,}336 ordinary working days,
1{,}216 marked U (the second weekend day and public holidays) and 1{,}087
marked A. This column is unusual and worth pausing on: \emph{we know it in
advance}. The next day's type is on the calendar today, so it is a legitimate
input to a forecast of that day. Most features are not like that.

Plotting twenty-one years shows three things, none of which is what we will
model: bus falling while rail rose and then fell, an annual wobble of about a
tenth of the level, and March 2020 as a \emph{cliff} rather than a dip. A model
fitted across that cliff is being asked to describe two different cities, so we
model 2016 to 31 May 2019 --- 1{,}247 days of one regime. Everything after is
set aside, not because it is uninteresting but because we cannot both use it
and claim to have been surprised by it. Data snooping applies to time as well
as to rows.

Zooming in until individual days are visible shows a seven-day cycle so strong
it dominates the plot, and one trough in the wrong place. Ask the data rather
than the calendar: the three days around it are marked A, U, U --- a public
holiday extending the weekend, which the column we were given already knows.
\emph{Every anomaly should end either in an explanation or in a note that you
could not find one. Not in silence.}

\section{Derivation: stationarity, differencing and autocorrelation}

Every model in this course assumes the data it was trained on and the data it
will meet are the same kind of thing. For a series, what does that assumption
say --- and does ours satisfy it?

\begin{definition}[Weak stationarity]
$\{y_t\}$ is weakly stationary if $\mathbb{E}[y_t] = \mu$ is constant,
$\operatorname{Var}(y_t) = \sigma^2$ is constant and finite, and
$\operatorname{Cov}(y_t, y_{t+k}) = \gamma_k$ depends on $k$ only.
\end{definition}

Only the first two moments, and only the requirement that the covariance
\emph{forgets when and remembers only how far apart}. Strict stationarity ---
the whole joint law invariant under a shift --- is far more than we need and
more than any real series has.

\begin{keybox}[title={Why a model depends on it}]
Fitting minimises an average loss over the training rows, and we then quote
that as an estimate of the loss on future rows. That step is a \emph{claim}:
the future rows come from the same distribution. For rows indexed by time,
``the same distribution'' is exactly what weak stationarity asserts about the
first two moments.

If $\mu$ drifts, a model fitted on the past is systematically wrong on the
future --- not noisily wrong, \emph{biased}. That is the difference between an
error that averages out and one that accumulates.
\end{keybox}

Our rolling mean moves by more than 170{,}000 riders across the record, so the
first condition fails outright. An augmented Dickey--Fuller test returns
$-4.61$ with $p \approx 0.0001$ and rejects its null --- but read what was
rejected: the null is the presence of a \emph{unit root}, one specific way of
being non-stationary. Rejecting it does not establish stationarity, which the
plot already refuted. A test answers the question it was given, and ours was
not that question.

\subsection{Differencing}

With $(\nabla y)_t = y_t - y_{t-1}$, applied to a sampled function $\nabla$ is
a discrete derivative. One value is lost at the start each time; count those
losses, because \code{diff} returns \code{NaN} there and says nothing.

\begin{derivbox}[title={$\nabla^{d}$ annihilates a polynomial trend of degree $d$}]
Monomials suffice, as $\nabla$ is linear:
\[
\begin{aligned}
t^{d} - (t-1)^{d}
 &= \sum_{j=0}^{d-1}\binom{d}{j}(-1)^{d-1-j}t^{\,j} \\
 &= d\,t^{\,d-1} + \text{lower order}.
\end{aligned}
\]
The leading coefficient is $d \neq 0$, so the degree drops by \emph{exactly}
one. Induct.
\end{derivbox}

A linear series becomes a constant, and the constant is the slope --- and a
constant mean is the first condition of weak stationarity, which we have just
bought. That $d$ is the \term{order of integration}, the I in ARIMA.

Seasonal differencing, $(\nabla_s y)_t = y_t - y_{t-s}$, removes a periodic
component \emph{exactly}: if $y_t = c_t + r_t$ with $c_t = c_{t-s}$, then
$(\nabla_s y)_t = r_t - r_{t-s}$, with no approximation and no fitted
parameter.

\begin{center}
\small
\begin{tabular}{lr}
\toprule
\textbf{Series} & \textbf{Standard deviation} \\
\midrule
the level & 183{,}841 \\
differenced once & 198{,}098 \\
differenced at lag 7 & 104{,}730 \\
\bottomrule
\end{tabular}
\end{center}

One difference \emph{increased} the variability by 8\%. That is not a bug and
not an accident --- there is an exact condition.

\subsection{Autocorrelation, and when differencing helps}

For a weakly stationary series,
$\rho_k = \gamma_k/\gamma_0 \in [-1,1]$ is the correlation of the series with
itself $k$ steps later. Note it is defined \emph{only} under stationarity:
without it the covariance depends on $t$ as well as $k$ and there is no single
number to plot.

\begin{derivbox}[title={The variance of a lag-$k$ difference}]
\[
\begin{aligned}
\operatorname{Var}(y_t - y_{t-k})
 &= \operatorname{Var}(y_t) + \operatorname{Var}(y_{t-k})
   - 2\operatorname{Cov}(y_t, y_{t-k}) \\
 &= 2\sigma^{2}\bigl(1 - \rho_k\bigr).
\end{aligned}
\]
So differencing at lag $k$ reduces the variance if and only if
$\rho_k > \tfrac12$.
\end{derivbox}

\begin{center}
\small
\begin{tabular}{rrrr}
\toprule
\textbf{Lag $k$} & $\rho_k$ & \textbf{Predicted sd} & \textbf{Measured} \\
\midrule
1 & 0.419 & 198{,}236 & 198{,}098 \\
7 & 0.831 & 106{,}782 & 104{,}730 \\
\bottomrule
\end{tabular}
\end{center}

The lag-1 prediction is right to one part in a thousand and the lag-7 to two
per cent, the residual discrepancy being the series failing to be exactly
stationary --- which we already knew. And $\rho_1 < \tfrac12$ while
$\rho_7 > \tfrac12$, so the first difference \emph{must} be worse and the
seventh \emph{must} be better. Both were, before we knew why.

\begin{keybox}[title={And now the naive forecast is explained}]
``Copy the value from seven days ago'' has error
$\hat y_t - y_t = -(\nabla_7 y)_t$ --- its error \emph{is} the seasonal
difference. So its error variance is $2\sigma^2(1-\rho_7)$, which at
$\rho_7 = 0.831$ is $0.34\,\sigma^2$.

The naive forecast is hard to beat for one reason and one reason only:
$\rho_7$ is large. It is not a lucky heuristic --- it is the autocorrelation
function, read off at a single lag.
\end{keybox}

\begin{failbox}[title={Failure condition --- an MAE inferred from a Gaussian that is not there}]
A Gaussian of standard deviation 104{,}730 would have mean absolute value
$\sqrt{2/\pi}\,\sigma \approx 83{,}562$. Our measured MAE is 55{,}051.

The errors are not Gaussian: the excess kurtosis of the seven-day difference is
12 --- most days are almost exactly like the same day last week, and a handful
are catastrophically different. A variance is the wrong summary of that
distribution, and so is any conversion assuming normality.

Which is also the argument for MAE over RMSE, arriving from a second
direction.
\end{failbox}

\section{Baselines, and the metric}

We predict $\hat y_{t+1} = f(y_{t-55},\dots,y_t)$: given the last eight weeks,
the next day. MAE is the choice, because the units are \emph{passengers} and
the operations team can act on ``we were out by forty thousand boardings'';
because the worst days are holidays the calendar already tells us about, and we
do not want the metric dominated by days you can look up; and because RMSE
would make one bad holiday worth more than a whole ordinary month. MAPE is
reported beside it, never as the objective --- it divides by the truth, so it is
most sensitive exactly where the model is weakest, and a metric like that makes
every improvement look like noise.

Three baselines, all scored on exactly the rows the models will be scored on:
predict a constant, copy the day before, copy last week.

\begin{keybox}[title={The number to beat}]
Copying the value from seven days ago gives an MAE of \textbf{55{,}399} over
1{,}191 days --- 12.0\% of a typical day, and 2.4 times better than copying the
previous day. The seven-day cycle is worth more than recency.

It needs no training, no maintenance, no monitoring and no data scientist, and
it fails only on days that appear on a public calendar. A transit authority
could put it into service in an afternoon.
\end{keybox}

One holiday costs it 464{,}640 in a single day --- more than eight ordinary
days of error put together. Any model has to be better than this \emph{by
enough to justify existing}, and you should write down what ``enough'' is
before you see any model output.

\subsection{Turning one series into a table}

The dataset class is nine lines, and the subtlety is entirely in
\code{\_\_len\_\_} and in the single index that is both the end of the window
and the position of the label. Test it on a series whose answer you know: a toy
of $0,1,2,3,4,5$ with window 3 gives $[0,1,2] \to 3$, $[1,2,3] \to 4$,
$[2,3,4] \to 5$ --- three rows, three targets, none inside its own window.
Twenty seconds, and it is the only thing standing between you and an off-by-one
that produces a beautiful score.

Eight weeks is a whole number of seven-day cycles, so every position sees the
same phase of the week; long enough to contain a slow change in level; and
short enough that 1{,}247 days still yield 1{,}191 rows. Longer windows cost
rows --- a 365-day window gives 882 instead of 1{,}191, so the model has more
to learn from less.

Dividing by a million puts every value near $[0,1]$ where the default
initialisation and learning rate expect their inputs. That is scaling, and
unlike \code{StandardScaler} it learns nothing from the data, so it cannot
leak.

\section{The protocol, and what it is worth}

$\rho_7 = 0.831$ says any two days a week apart are strongly dependent.
Cross-validation assumes the held-out rows are independent of the rows the
model was fitted on. \emph{Those two sentences cannot both be acted on.}

\begin{failbox}[title={Failure condition --- a rule applied where its reason does not hold}]
Lecture 2's checklist says: ``Seeds for every generator;
\code{KFold(shuffle=True, random\_state=42)}, never the default.''

That rule was correct, and it solved a real problem: a frame sorted differently
gave different folds and nobody could see why. \emph{Here it is the bug.}

A rule you cannot say the reason for is a rule you will apply where it does not
hold. This is the clearest example the course contains.
\end{failbox}

Cross-validation is unbiased for the risk under two conditions: the held-out
rows are independent of the training rows, and they are drawn from the
distribution the model will meet. A shuffled split of a time series satisfies
neither.

Condition 1 fails and $\rho_k$ says by how much: the covariance is
$0.83\sigma^2$ at lag 7 and never near zero for any lag inside the window. A
model fitted on rows \emph{surrounding} a held-out row faces a smaller
conditional variance than one fitted only on rows before it, so the estimate is
biased downwards --- optimistic in the overwhelming majority of cases. Not
``always'': that would be a theorem, and the conditional-variance argument is a
heuristic rather than a bound.

Condition 2 fails and it is not subtle. In deployment every training day
precedes every day being forecast; under a shuffled split, four training days
in five come \emph{after} the day being scored. The model is not forecasting.
It is being handed the shape of the series around the point it is asked about,
and interpolating --- a different and much easier problem, and nobody is paying
for it.

\begin{center}
\small
\begin{tabular}{lrrr}
\toprule
\textbf{Model} & \textbf{Random split} & \textbf{Forecasting forward} &
\textbf{Gap} \\
\midrule
Linear, 56 lags & 44{,}925 & 56{,}446 & $+25.6\%$ \\
Simple RNN, 32 units & 38{,}224 & 49{,}073 & $+28.4\%$ \\
\bottomrule
\end{tabular}
\end{center}

\subsection{Now do the thing this course exists for}

We have a gap of 28\%. \emph{How much of it is the split?} Three other things
changed when the protocol changed, and each moves the number on its own.

\textbf{Confound 1 --- the test rows are harder.} Copying last week needs no
fitting, so its error measures only how difficult a set of days is. On the
shuffled folds it scores 55{,}385; on the five months after the cut, 64{,}815.
The forecast period is 17\% harder before any model is involved, because it
contains the winter holidays. This is Lecture 2's ``we dropped the capped
districts'' returning: any time a number moves, ask first whether the
evaluation set moved.

\textbf{Confound 2 --- the training set changed size.}
\code{TimeSeriesSplit}'s first fold fits on a fifth of the rows and scores
80{,}506, which is not a statement about the protocol but about having 238
training rows. The repair is a rolling-origin backtest with a \emph{fixed}
training window --- 700 rows for every origin, five origins, 98 days scored
each time --- so the folds differ only in \emph{when} they are.

\textbf{Confound 3 --- the folds have a spread.} The random folds span 39{,}538
to 51{,}096. A difference in means smaller than that spread is not a
difference: the paired-comparison lesson from Lecture 2, unchanged.

\begin{center}
\small
\begin{tabular}{lrrr}
\toprule
\textbf{Protocol} & \textbf{Model} & \textbf{Naive} & \textbf{Ratio} \\
\midrule
random draws of 98 rows, 20 of them & 48{,}093 & 61{,}461 & 0.782 \\
rolling origin, 5 origins & 47{,}212 & 55{,}371 & 0.853 \\
\bottomrule
\end{tabular}
\end{center}

Same training size, same test size, same model --- only the arrangement differs
--- and each scored against copying last week \emph{on its own rows}, which
cancels the difficulty of the period. The random split reports the model as
9.0\% more skilful than it is.

\begin{keybox}[title={The honest verdict, in three numbers}]
\textbf{28\%} is the raw gap on the number you committed. That is what you
would have been wrong by, and it is real.

\textbf{17\%} of it is the forecast period genuinely being harder --- the
baseline moves too.

\textbf{9 to 10\%} is the protocol itself, and it survives every control we
could apply: matched training size, matched test size, and the baseline
computed on identical rows. The linear model is overstated by 7.4\% and the
recurrent network by 9.7\%, which is what you expect if the cause is the split
rather than the model.
\end{keybox}

Put in the terms a stakeholder hears: we told them the recurrent model was
31\% better than doing nothing, and it is 24\% better. A fifth of the case for
building it was an artefact of the evaluation, and our own linear model loses a
third the same way, 18.9\% to 12.9\%.

\begin{center}
\small
\begin{tabular}{p{0.22\textwidth}p{0.30\textwidth}p{0.34\textwidth}}
\toprule
& \textbf{Scaling before the split} & \textbf{Random split on a series} \\
\midrule
Diagnosed in & Lecture 2 & here \\
Cost, measured & nothing a measurement could find & a fifth to a third of
                                                    everything a model was
                                                    worth \\
Why that size & an invertible affine map on an equivariant estimator &
                $\rho_k$ large at every lag inside the window \\
\bottomrule
\end{tabular}
\end{center}

Same procedural rule, two wildly different costs. \emph{You cannot tell which
case you are in from the number --- only from the structure. That is why the
rule is procedural.}

\section{The linear model}

\[ \hat y_{t+1} = b + \sum_{k=0}^{55} w_k\,y_{t-k}. \]
Fifty-six weights and a bias --- and if the model learns $w_6 = 1$ with every
other weight zero, it has reinvented copying last week. \emph{The linear model
contains the baseline}, so it cannot do worse unless we fit it badly.

It scores 44{,}925 against 55{,}399: 18.9\% better, with folds from 39{,}538 to
51{,}096, and over twenty different shuffles the mean moves by about 300 --- so
the improvement is far larger than the noise in the estimate.

What did the weights learn? The largest is at lag 1 at $+0.45$, so the day
before still carries the level. After that the next four largest are at lags
35, 7, 14 and 28 --- all seven multiples of seven have a positive weight and no
other lag does so consistently. The weights sum to 0.86 with an intercept
around 90{,}000 riders, so the model shrinks gently towards the mean.

\emph{It rediscovered the seven-day cycle from the data. Nobody told it that a
week has seven days.}

\section{The notebook}

\begin{trybox}[title={What to change}]
Shuffle the rows before splitting and re-run the evaluation. The score
improves, and nothing warns you. That improvement is the whole argument of the
second half of this lecture, produced by you rather than quoted at you.
\end{trybox}

\section{Exercises}

\begin{enumerate}[leftmargin=1.6em,itemsep=4pt]
  \item State the three conditions of weak stationarity, and explain why a
        model that quotes a training average as an estimate of future loss
        depends on them. \hfill \scope{[5]}
  \item Prove that $\nabla$ reduces the degree of a polynomial by exactly one,
        and state what $\nabla^d$ therefore does to a degree-$d$ trend.
        \hfill \scope{[4]}
  \item Derive $\operatorname{Var}(y_t - y_{t-k})$ in terms of $\rho_k$, and
        give the exact condition under which differencing at lag $k$ reduces
        the variance. \hfill \scope{[5]}
  \item A colleague evaluates a forecasting model with
        \code{KFold(shuffle=True)} and reports a 31\% improvement over a naive
        baseline. Name the two conditions the protocol violates, and say in
        which direction the estimate is biased. \hfill \scope{[4]}
  \item A model scored under a forward protocol is 28\% worse than under a
        random split. Name three things other than the protocol that could
        explain part of that gap, and how you would control for each.
        \hfill \scope{[4]}
\end{enumerate}

\section*{Summary}
\addcontentsline{toc}{section}{Summary}

Weak stationarity asks for a constant mean, a constant finite variance and a
covariance depending only on the lag --- and models depend on it because they
quote a training average as an estimate of a future one. Our series does not
have it. $\nabla^d$ annihilates a polynomial trend of degree $d$ because
$\nabla$ drops the degree by exactly one, and $\nabla_s$ annihilates a
component of period $s$ exactly.

$\operatorname{Var}(y_t - y_{t-k}) = 2\sigma^2(1-\rho_k)$, so differencing
helps if and only if $\rho_k > \tfrac12$ --- which predicted, before it was
measured, that the first difference would be worse and the seventh better. The
naive forecast's error \emph{is} the seasonal difference, so its strength is
$\rho_7 = 0.831$ and nothing else.

Then the lecture's real subject. Lecture 2's shuffled-fold rule, correct where
it was written, is the bug here: it violates both conditions cross-validation
needs, and it reported the recurrent model as 31\% better than nothing when it
is 24\% better. Of the raw 28\% gap, 17\% is the forecast period genuinely
being harder and 9 to 10\% is the protocol, surviving matched training size,
matched test size and a baseline recomputed on identical rows.

\end{document}
